% Options for packages loaded elsewhere
\PassOptionsToPackage{unicode}{hyperref}
\PassOptionsToPackage{hyphens}{url}
\PassOptionsToPackage{dvipsnames,svgnames,x11names}{xcolor}
%
\documentclass[
  letterpaper,
  DIV=11,
  numbers=noendperiod]{scrartcl}

\usepackage{amsmath,amssymb}
\usepackage{iftex}
\ifPDFTeX
  \usepackage[T1]{fontenc}
  \usepackage[utf8]{inputenc}
  \usepackage{textcomp} % provide euro and other symbols
\else % if luatex or xetex
  \usepackage{unicode-math}
  \defaultfontfeatures{Scale=MatchLowercase}
  \defaultfontfeatures[\rmfamily]{Ligatures=TeX,Scale=1}
\fi
\usepackage{lmodern}
\ifPDFTeX\else  
    % xetex/luatex font selection
\fi
% Use upquote if available, for straight quotes in verbatim environments
\IfFileExists{upquote.sty}{\usepackage{upquote}}{}
\IfFileExists{microtype.sty}{% use microtype if available
  \usepackage[]{microtype}
  \UseMicrotypeSet[protrusion]{basicmath} % disable protrusion for tt fonts
}{}
\makeatletter
\@ifundefined{KOMAClassName}{% if non-KOMA class
  \IfFileExists{parskip.sty}{%
    \usepackage{parskip}
  }{% else
    \setlength{\parindent}{0pt}
    \setlength{\parskip}{6pt plus 2pt minus 1pt}}
}{% if KOMA class
  \KOMAoptions{parskip=half}}
\makeatother
\usepackage{xcolor}
\setlength{\emergencystretch}{3em} % prevent overfull lines
\setcounter{secnumdepth}{-\maxdimen} % remove section numbering
% Make \paragraph and \subparagraph free-standing
\makeatletter
\ifx\paragraph\undefined\else
  \let\oldparagraph\paragraph
  \renewcommand{\paragraph}{
    \@ifstar
      \xxxParagraphStar
      \xxxParagraphNoStar
  }
  \newcommand{\xxxParagraphStar}[1]{\oldparagraph*{#1}\mbox{}}
  \newcommand{\xxxParagraphNoStar}[1]{\oldparagraph{#1}\mbox{}}
\fi
\ifx\subparagraph\undefined\else
  \let\oldsubparagraph\subparagraph
  \renewcommand{\subparagraph}{
    \@ifstar
      \xxxSubParagraphStar
      \xxxSubParagraphNoStar
  }
  \newcommand{\xxxSubParagraphStar}[1]{\oldsubparagraph*{#1}\mbox{}}
  \newcommand{\xxxSubParagraphNoStar}[1]{\oldsubparagraph{#1}\mbox{}}
\fi
\makeatother


\providecommand{\tightlist}{%
  \setlength{\itemsep}{0pt}\setlength{\parskip}{0pt}}\usepackage{longtable,booktabs,array}
\usepackage{calc} % for calculating minipage widths
% Correct order of tables after \paragraph or \subparagraph
\usepackage{etoolbox}
\makeatletter
\patchcmd\longtable{\par}{\if@noskipsec\mbox{}\fi\par}{}{}
\makeatother
% Allow footnotes in longtable head/foot
\IfFileExists{footnotehyper.sty}{\usepackage{footnotehyper}}{\usepackage{footnote}}
\makesavenoteenv{longtable}
\usepackage{graphicx}
\makeatletter
\def\maxwidth{\ifdim\Gin@nat@width>\linewidth\linewidth\else\Gin@nat@width\fi}
\def\maxheight{\ifdim\Gin@nat@height>\textheight\textheight\else\Gin@nat@height\fi}
\makeatother
% Scale images if necessary, so that they will not overflow the page
% margins by default, and it is still possible to overwrite the defaults
% using explicit options in \includegraphics[width, height, ...]{}
\setkeys{Gin}{width=\maxwidth,height=\maxheight,keepaspectratio}
% Set default figure placement to htbp
\makeatletter
\def\fps@figure{htbp}
\makeatother

\KOMAoption{captions}{tableheading}
\makeatletter
\@ifpackageloaded{caption}{}{\usepackage{caption}}
\AtBeginDocument{%
\ifdefined\contentsname
  \renewcommand*\contentsname{Table of contents}
\else
  \newcommand\contentsname{Table of contents}
\fi
\ifdefined\listfigurename
  \renewcommand*\listfigurename{List of Figures}
\else
  \newcommand\listfigurename{List of Figures}
\fi
\ifdefined\listtablename
  \renewcommand*\listtablename{List of Tables}
\else
  \newcommand\listtablename{List of Tables}
\fi
\ifdefined\figurename
  \renewcommand*\figurename{Figure}
\else
  \newcommand\figurename{Figure}
\fi
\ifdefined\tablename
  \renewcommand*\tablename{Table}
\else
  \newcommand\tablename{Table}
\fi
}
\@ifpackageloaded{float}{}{\usepackage{float}}
\floatstyle{ruled}
\@ifundefined{c@chapter}{\newfloat{codelisting}{h}{lop}}{\newfloat{codelisting}{h}{lop}[chapter]}
\floatname{codelisting}{Listing}
\newcommand*\listoflistings{\listof{codelisting}{List of Listings}}
\makeatother
\makeatletter
\makeatother
\makeatletter
\@ifpackageloaded{caption}{}{\usepackage{caption}}
\@ifpackageloaded{subcaption}{}{\usepackage{subcaption}}
\makeatother

\ifLuaTeX
  \usepackage{selnolig}  % disable illegal ligatures
\fi
\usepackage{bookmark}

\IfFileExists{xurl.sty}{\usepackage{xurl}}{} % add URL line breaks if available
\urlstyle{same} % disable monospaced font for URLs
\hypersetup{
  pdftitle={Assembly and Annotation of a Leptospira Isolate from a New Hampshire Fisher (Pekania pennanti)},
  colorlinks=true,
  linkcolor={blue},
  filecolor={Maroon},
  citecolor={Blue},
  urlcolor={Blue},
  pdfcreator={LaTeX via pandoc}}


\title{Assembly and Annotation of a Leptospira Isolate from a New
Hampshire Fisher (Pekania pennanti)}
\author{}
\date{}

\begin{document}
\maketitle


\section{Assembly and Annotation of a Leptospira Isolate from a New
Hampshire Fisher (Pekania
pennanti)}\label{assembly-and-annotation-of-a-leptospira-isolate-from-a-new-hampshire-fisher-pekania-pennanti}

\subsection{Author}\label{author}

Megan S. Munis, DVM MS. University of New Hampshire College of Life
Sciences and Agriculture.

\subsection{Introduction}\label{introduction}

Leptospira is a genus of zoonotic spirochete bacteria that includes
several species capable of causing disease in humans and animals. This
study focuses on the assembly and annotation of a Leptospira isolate
obtained from a New Hampshire Fisher (Pekania pennanti). Disseminated in
the urine of infected humans and animals, innoculation of the mucus
membranes results in systemic infection and a range of symptoms, from
mild flu symptoms to severe illness, renal and hepatic failure, and
death. Historically mesocarnivoies such as fisher have been suspected to
be chronic carriers of Leptospira and consequently serve as reservoirs
for the pathogen in the environment and a source of infection for humans
and animals, however previous studies are limited to serologic or
molecular techniques which show exposure to the bacteria and its
presence in the animals respectively, but do not explain the
pathogenesis of the bacteria and its potential to cause disease in
wildlife species. By culturing leptospira from a live fisher we have the
unique opportunity to study the genomic characteristics of a Leptospira
isolate from wildlife and understand their pathogenic potential.
Understanding the genomic characteristics of Leptospira isolates from
wildlife can provide insights into their pathogenic potential and inform
public health strategies for managing leptospirosis. The objectives of
this study were to assemble the genome of the Leptospira isolate, assess
its completeness, and annotate its gene content to better understand its
potential virulence factors and ecological role.

\subsection{Methods}\label{methods}

Sample Collection and Bacterial Culture Urine was collected by sterile
cystocentesis from a live-trapped fisher in Durham, NH as a component of
a collaborative study between New Hampshire Fisher and Game and the
University of New Hampshire. Urine was inoculated into
Ellinghausen-McCullough-Johnson-Harris (EMJH) media and incubated at
30°C for 4 weeks at USDA-ARS National Animal Disease Center (NADC) in
Ames, IA. The resulting cultured isolate was frozen at -80°C and shipped
to the University of New Hampshire for DNA extraction and sequencing.

Sequencing The DNA extraction and sequencing were performed at the
University of New Hampshire Hubbard Center for Genome Studies. The
isolate was sequenced using an Illumina HiSeq 2500 platform with 250p
paired-end reads. The raw sequencing data were processed using
Trimmomatic for quality control and adapter trimming. The cleaned reads
were then assembled de novo using SPAdes assembler. Assembly quality was
assessed using BUSCO and QUAST tools. The resulting contigs were
annotated using Prokka, to provide gene predictions and functional
annotations. A BLAST query was completed against the 16s rRNA gene to
confirm the taxonomic identity of the isolate. *filtered samtools to
500bp, depth of 10.

\subsection{Results}\label{results}

QUAST results revealed 87 contigs with an N50 of 122,144 bp, and a L50
of 11bp. The GC content of 35.92\%, and total assembly size of 4.36bp is
consistent with other Leptospira genomes (Mejia 2022). BUSCO analysis
revealed 95.2\% of the expected single-copy orthologs were present in
the assembly, indicating a high level of completeness. Prokka annotation
identified 3,500 protein-coding genes, 50 tRNA genes, and 5 rRNA
operons. The BLAST query of the 16s rRNA gene confirmed that the isolate
belongs to the Leptospira genus, with a 99\% identity match to
Leptospira kirschneri.

\includegraphics{plots/fst_div.png} Figure 1. Provide figure
descriptions- indicate how figure was made (programs), with which data.
Make sure to cite everything you used.

\subsubsection{Note:}\label{note}

To add a plot from the `figs' folder, there needs to be a line in this
document that looks like this:

\begin{verbatim}
![plot](figs/plotfile.png)
\end{verbatim}




\end{document}
